\chapter*{Preface}
\addcontentsline{toc}{chapter}{Preface}
\markboth{Preface}{Preface}

A topology optimization can converge and still solve the wrong problem. A support may restrain a direction the
engineer meant to leave free, a load may sit on the wrong face, a pad may be retained that nobody asked for. Neither the
solver's residual nor a language agent's confident summary reveals such errors. AgenticTO is built around that
observation. It keeps the design intent in one typed specification, checks it with deterministic rules before anything
is solved, and lets only the solver issue the certificate that marks a result as released. A language agent can help
at every step, but it works through the same logged tools as the engineer and cannot write a certificate.

This guidebook explains how to install the software, build and solve designs in its graphical workspace, work with the
built-in assistant, connect other agents through the Model Context Protocol, read the provenance log, and reproduce the
results of the companion paper.

\paragraph{Who this book is for.} Two groups of readers. Journal reviewers who want to know what the software actually
does, which check is responsible for which guarantee, and how each number in the paper can be traced to a file and a
script: Chapters~\ref{ch:overview}, \ref{ch:ladder}, \ref{ch:reproduce} and \ref{ch:limits} are written for you.
Engineers who download the public release and want to use it: start with Chapters~\ref{ch:install} and
\ref{ch:quickstart}, then read the chapters on the workspace, the specification and the rules.

\paragraph{How to read it.} Chapter~\ref{ch:overview} gives the concepts in a few pages. Chapter~\ref{ch:quickstart}
takes you from a fresh start to a certified, exported design in about ten minutes. The remaining chapters can be read in
any order. The appendices are reference material: every field of the specification, every tool with its arguments,
every certificate field, the keyboard map, every script, and a glossary. A one-page quick reference precedes the
index.

\paragraph{Conventions.}
\begin{itemize}
  \item Interface labels are set in bold sans serif: \gui{Optimize release}. Keys look like \keycap{S} or
        \keycap{Ctrl}+\keycap{K}.
  \item File names, commands, tool names and JSON fields are set in monospace: \code{scripts/serve.py},
        \tool{validate_spec}, \code{volume\_fraction}. Tool names link to their reference entry in
        Appendix~\ref{app:tools}; rule codes such as \rref{R003} link to Chapter~\ref{ch:rules}.
  \item Commands are shown for a POSIX shell. On Windows, the same commands work in PowerShell except where a box says
        otherwise.
  \item A grey bracket such as \claim{C2.3} names the claim in \code{results/claims\_ledger.md} that backs a number.
\end{itemize}

\begin{tipbox}[Four kinds of boxes]
Green boxes like this one give practical tips. Blue boxes add background. Vermillion boxes warn about something that
can cost you a result or mislead you. Amber boxes quote evidence from the paper, with its denominator.
\end{tipbox}

\paragraph{Acknowledgement of tools.} The software was developed with the help of OpenAI Codex and Claude-based
coding tools, as stated in the paper's declaration. This guidebook was drafted with AI assistance from the source code, the
existing user guide and the paper; its reference tables are generated from the code.
